% =============================================================================
% Foreword to Computational Thermodynamics --- Body and Soul Vol 5
% =============================================================================
% Positions the volume as the macroscopic-continuum companion to RealQM (Vol 6):
% a deterministic continuum reformulation of thermodynamics in which the 2nd
% Law is derived from finite-precision computation rather than postulated from
% statistical mechanics, presented as the realisation of the Body and Soul
% programme made tractable by AI collaboration and the GPU revolution.
% =============================================================================

\chapter*{Foreword}
\addcontentsline{toc}{chapter}{Foreword}
\markboth{FOREWORD}{FOREWORD}

% Opening 4 paragraphs (What ThermoDynamics is / Turbulent dissipation /
% Partial irreversibility / Physical computation, digital simulation)
% promoted to Chapter 1 "Essence of Real Thermodynamics" in chapters/essence.tex.

The Body and Soul programme, begun three decades ago, has from the start been
organised around a single methodological idea: to combine the \emph{Soul} of
analytical mathematics with the \emph{Body} of computational mathematics, so
that every part of mathematical physics is presented in a form which pairs the
underlying analytical structure with a runnable computational realisation.
Volumes I--III carried this commitment through \emph{Calculus} and
\emph{Linear Algebra}, each subject presented in both its analytical and its
computational form. \textbf{Volume IV}, on \emph{Turbulent Incompressible
Flow}~\cite{ambsflow}, took the programme one step further: the computational
form was put forward as the \emph{true mathematical model of physics}, based
on the view that physical reality, evolving in time, itself performs an
analog computation which the digital computer can parallel or simulate.
Turbulence, which had resisted understanding within the purely analytical
framework, could then be understood and simulated through a computational
form of the Navier--Stokes equations. The \textbf{present Volume V} extends
the same mathematics from incompressible to compressible flow, where it
becomes the essence of \emph{thermodynamics} --- with an understandable
computational form of the 2nd Law as the inevitable dissipation of
finite-precision computation. \textbf{Volume VI}, \emph{Real Quantum
Mechanics}~\cite{RealQM2026}, extends the same mathematics again, from
macroscopic compressible flow down to the atomic scale, replacing the
classical wave function in $3N$-dimensional configuration space with a
multiphase 3D continuum formulation in real three-dimensional space. The
two volumes V and VI together complete the Body and Soul programme by
carrying the combined analytical/computational point of view across the
full scale range from macroscopic flows to atomic structure, with
statistical aggregation replaced by deterministic field equations at both
ends.

\paragraph{The thesis in one sentence.}
The 2nd Law of thermodynamics, irreversibility, and the arrow of time are not consequences of statistical aggregation over configurations of atoms. They are consequences of \textbf{finite-precision computation} on deterministic continuum field equations: any finite-resolution numerical solution of the Euler equations of compressible flow dissipates kinetic energy into internal energy through the small-scale residual that the resolution cannot represent, and this dissipation is irreversible by the structure of the discretisation, not by any statistical-mechanical assumption. The book argues that this is the right reading of the 2nd Law: a structural consequence of the finite information content of any computable representation of the physical world, including the physical world's own representation of itself as analog computation.

\paragraph{The two central problems: Joule's experiment and the piston-cylinder.}
The book is organised around two foundational computational problems that
together carry the conceptual content of thermodynamics:

\begin{itemize}
\item \textbf{Joule's 1845 experiment} (Chapter~\ref{ch:joule}) — gas in a
high-pressure chamber expands through a channel into a low-pressure chamber.
Bulk flow through the channel converts internal energy into kinetic energy
during the expansion; turbulent/shock dissipation in the receiving chamber
converts it back into heat during the recompression. The two halves are
asymmetric in the dynamic dissipation pattern, and the result is a residual
temperature gradient between the chambers that classical equilibrium
thermodynamics cannot predict. This is the \textbf{refrigerator principle}:
the same expansion-cooling and compression-heating asymmetry that a
refrigeration cycle exploits with external work input, demonstrated in
one-shot form by simply releasing a pressure differential between two
chambers.

\item \textbf{The piston-cylinder} (Chapter~\ref{ch:piston}) — gas in a chamber
pushes a movable wall (piston) outward against an external load, doing work
on the load. With a smooth single-chamber geometry the gas does work
quasi-statically and the process approaches reversible efficiency; with a
channel constriction between the gas reservoir and the piston (the
Joule geometry with a movable far wall), the gas must jet through the
constriction and recompress against the piston, with turbulent
dissipation pulling the efficiency $\eta = W/\Delta U$ below the
reversible limit. This is the \textbf{heat-engine principle}: the
conversion of stored internal energy into useful work delivered to an
external load, with irreversibility entering through the same
turbulent-dissipation channel as in Joule's experiment.
\end{itemize}

The two problems are duals of each other in the deterministic-continuum
reformulation. The Joule experiment shows what happens when the
asymmetric expansion-compression cycle is allowed to dissipate freely; the
piston-cylinder shows what happens when the same dynamics are channelled
into useful work against a load. The 2nd Law in its RNS reading — finite-
precision dissipation in the discretisation — sets the irreversibility in
both cases. Refrigeration cycles run the first problem cyclically with
external work input; heat engines run the second problem cyclically with
heat input from a reservoir. Vol VI Applications develops both into full
cycles (Carnot heat engines, Carnot heat pumps, compressor refrigerators);
the present opening chapters establish the foundational single-shot dynamics
that the cycles are built from.

\paragraph{Real Thermodynamics in the Body and Soul tradition.}
The earlier volumes of the programme presented continuum mechanics for macroscopic matter: fluids, elastic solids, plasticity, transport. This volume presents thermodynamics inside the same continuum-mechanics framework, without adding a statistical-mechanical layer. The Euler and Navier--Stokes equations carry the dynamics; the \emph{Real Navier--Stokes} method (RNS) --- `Real' here meaning the computational realisation as a least-squares stabilised Galerkin method --- supplies the finite-precision computation; the 2nd Law emerges as the entropy condition for weak solutions in the same sense that it does for shocks in classical conservation-law theory. The Boltzmann--Gibbs apparatus --- microstates, ensembles, the probabilistic foundation of equilibrium thermodynamics --- is not invoked at the foundational level. Where statistical reasoning is the natural framing for a specific regime (large-system equilibrium fluctuations, critical phenomena) it remains available; what is replaced is the claim that statistics is the foundation of the discipline.

\paragraph{The programme as mathematics education, and as it is read today.}
The Body and Soul mathematics education programme was initiated thirty years ago, around the proposition that a modern undergraduate education in mathematics should be a synthesis of the formal analytical mathematics of Calculus --- the \emph{Soul} of the subject --- with the computational mathematics made possible by the computer --- the \emph{Body}. The programme was met with resistance from both directions. Established mathematics teachers, trained in formal analysis but not in programming, found the computational half foreign to their craft and slow to incorporate. Computer scientists, comfortable with programming but uncomfortable with the apparatus of Calculus, regarded the analytical half as an obstacle to be routed around. The synthesis languished in the gap between the two cultures it had been designed to bridge. The idea was, however, correct, and with the recent advent of capable AI assistance it acquires a new and stronger sense: the analytical and the computational sides of mathematics no longer have to be carried by the same human in the same head. The mathematician's mind frames the problem, formulates the equations, identifies what must be shown; the AI translates the formulation into running code, executes the calculations, and brings the construction alive on the screen. The synthesis that the Body and Soul programme proposed in education thirty years ago becomes, with AI as collaborator, a practical working method for research.

\paragraph{The Body side and the FEniCS project.}
The Body side of the programme --- the commitment that the computational realisation should be an organic part of the mathematics, not a separate engineering exercise --- has its principal historical embodiment in the \textbf{FEniCS Project}~\cite{FEniCS2012}, an open-source platform for the automated solution of partial differential equations by the finite element method. FEniCS was developed by the author's research group over roughly a decade of dedicated work, with the goal that a researcher who could write down a variational form on paper could obtain a running finite-element solver from it directly. The RNS method that supplies the computational backbone of this book was developed in that environment, and the original \emph{Computational Thermodynamics} (Hoffman, Johnson and Nazarov, 2012) on which the present volume is based reported numerical results obtained with the FEniCS-Unicorn implementation.

The relevance to the present rewrite is the change of timescale. As Anders Logg, principal developer of FEniCS, has observed, what required a decade of focused effort in the 2000s can today, with the help of capable AI assistance, be carried out in a day. The browser-based interactive simulators that accompany this volume in the public Gallery --- compressible flow demonstrations of Joule's 1845 experiment, of bluff-body Mach-2 turbulence, and of the Burgers equation as the canonical model PDE --- were produced on that compressed timescale. The Body side of the programme has not changed; what has changed is how much of it a single mathematical mind, with AI as collaborator, can carry from formulation to running code in a working session.

\paragraph{The GPU revolution: the third enabling condition.}
The Mind--AI collaboration is one of the enabling conditions for the present volume; the other is the GPU revolution of the past two decades. Graphics Processing Units originated as dedicated accelerators for 3D-graphics rendering in the 1990s and were repurposed for general scientific computation in the mid-2000s (CUDA 2007, OpenCL shortly after). The architectural feature that distinguishes a GPU from a conventional CPU is massive parallelism: thousands of identical compute threads executing the same instruction stream on different data, with memory bandwidth tuned for streaming through large regular arrays. The compressible Euler and Navier--Stokes equations on a regular grid are the canonical workload the GPU is built for. The 2023 finalisation of \textbf{WebGPU} --- a cross-vendor standard that exposes GPU compute through the browser without driver installation or local compilation --- makes it possible to ship a chapter's worth of computational thermodynamics as an HTML page the reader opens and runs on their own hardware at interactive rates, with no install step, no virtual environment, no submission to a remote service. The Joule's-experiment and bluff-body-Mach-2 simulators listed in Appendix~A run in this form. The reformulation makes the thermodynamics computable in principle; the GPU makes it computable on a laptop; AI assistance makes the laptop implementation feasible for one author in a working session. The three conditions together are what makes a Body-and-Soul-style thermodynamics volume possible at this moment.

\paragraph{Leibniz's project, realised.}
The proposition that the labour of calculation should be lifted from the mathematician's shoulders by a suitably constructed machine is, in its modern form, three and a half centuries old. Gottfried Wilhelm Leibniz, in the 1685 explanation of his Stepped Reckoner, wrote that ``it is unworthy of excellent men to lose hours like slaves in the labour of calculation which could safely be relegated to anyone else if machines were used.'' His wider programme --- the \emph{calculus ratiocinator}, a mechanical calculus for deriving conclusions from a \emph{characteristica universalis}, a universal symbolic language for reasoning --- proposed not merely arithmetic automation but the automation of mathematical thought as a whole, with the human mind freed to set the scene and the machine to carry the calculation. The intellectual lineage from this proposal through Babbage, Boole, Frege, Hilbert, G\"odel, and Turing to modern electronic and AI-assisted computation is the subject of Martin Davis's \emph{The Universal Computer: The Road from Leibniz to Turing}~\cite{Davis2000}. The collaboration between the human mathematical mind and Claude as AI assistant that produced the present volume is, in this lineage, Leibniz's three-hundred-and-forty-year-old proposal realised --- not as a special device for arithmetic, but as a working partnership in which the mind frames the equations and the machine, increasingly fluent in the language of mathematics, brings them alive.

\paragraph{Schr\"odinger's question, deferred.}
Schr\"odinger's \emph{What is Life?} (1944) framed living matter as
matter that sustains itself by feeding on \emph{negative entropy}: a
question that connects the 2nd Law as developed in this volume to the
cell biology taken up at the atomic scale in Volume~VI.

\paragraph{The synthesis that the Body and Soul programme completes.}
The earlier volumes of the programme presented continuum mechanics for macroscopic matter from the calculus side and from the algorithm side simultaneously. This volume does the same for thermodynamics, with the 2nd Law derived from the discretisation rather than postulated from statistical mechanics. Volume~VI does it for atomic-scale matter, with the wave function on physical $\mathbb{R}^3$ and the probability interpretation removed. The fluid in a pipe and the electron in an atom are, in the language the programme establishes across its volumes, the same kind of object: a field on $\mathbb{R}^3$ satisfying a partial differential equation, with boundary conditions that close the variational problem, and an algorithm that solves it on a computer. The programme that started with the Navier--Stokes equations and continued through classical elasticity now ends with the atomic Schr\"odinger equation in its real-3D-space form, with the deterministic continuum reformulation of thermodynamics in this volume as the bridge that connects the macroscopic and the atomic ends in a single methodological commitment.

\paragraph{The book contains no figures or detailed numerical results from the computations.}
All quantitative output --- velocity fields, density slices, energy-vs-time records, normal-mode animations, Joule expansion trajectories, bluff-body turbulence wakes --- is documented in detail in the public Gallery indexed in Appendix~A. The decision to keep the printed volume free of static figures is deliberate: static images of interactive simulations omit precisely what makes the framework distinctive. The chapter text describes what is computed and why it agrees with observation; the Gallery shows what it looks like, lets the reader vary parameters, and renders the framework's behaviour live. The two together constitute the book in its intended form.

\paragraph{A note on acknowledgements.}
The original \emph{Computational Thermodynamics} of 2012 was a joint work of Johan Hoffman, Claes Johnson, and Murtazo Nazarov; the present volume is its 2026 rewrite as Volume~V of the Body and Soul programme, by Claes Johnson in extended collaboration with \textbf{Claude}, the large-language-model AI assistant produced by Anthropic. The mathematical reformulation, the architecture, and the writing are the work of the author. The implementation of the browser-based interactive simulators, the systematic translation of the 2012 manuscript into the form the present volume takes, the cross-volume integration with Volume~VI, and a substantial fraction of the iterative back-and-forth that produced the rewrite were carried out with Claude. No human colleagues, students, funding agencies, departmental hosts, or institutional sponsors contributed to the rewrite in a way that would warrant the customary form of acknowledgement, beyond the recognition of the original 2012 co-authors recorded on the title page. The pattern of the human--AI collaboration follows the template documented in detail in the corresponding chapter of Volume~VI~\cite{RealQM2026}; that chapter is, in effect, the acknowledgements section of the present volume too.

\bigskip

\hfill \emph{Claes Johnson}\\
\hfill Stockholm, 2026

% --- End of Foreword ---
